\section{相关工作}

\subsection{在线跟踪中的关联与重连}

检测驱动的跟踪流水线由检测器和一条数据关联规则组成。SORT \cite{bewley2016sort} 定下现代形态：卡尔曼滤波预测每条轨迹，检测框按重叠度做匈牙利分配。DeepSORT \cite{wojke2017deepsort} 在同一级联里加入外观度量。ByteTrack \cite{zhang2022bytetrack} 打破了"低分检测即丢弃"的规则，同时引入让重连成为常态的轨迹缓冲：一条轨迹在消失后存活三十帧，之后任何越过关联阈值的检测都会被接回去。BoT-SORT \cite{aharon2022botsort} 改进了运动模型并补偿相机运动；OC-SORT \cite{cao2023ocsort} 用观测中心的估计替代滤波状态参与匹配。

这些系统改变的是"什么算一次好匹配"。重连是这一选择的副产品：跟踪器让丢失的轨迹继续存活，任何越过阈值的检测都被接受。缓冲长度和阈值决定了这种情况发生的频率，但两者都没有回答"这一次回来的是不是同一个人"。这就是本文的切入点。

\subsection{基于外观的重识别}

面对有歧义的回归，标准答案是外观描述子。DeepSORT \cite{wojke2017deepsort} 提出这一思路，此后的跟踪器或在运动之外保留一条重识别分支，或在跟踪结束后用学习到的相似度把轨迹片段链接起来。在具备训练好的嵌入模型、且两端裁剪都有足够像素时，这个信号很强，能处理几何无法解决的遮挡缺口。

它同时带来我们不作假设的前提：需要网络、需要特征库或重识别预算，还需要缺口两端都有可用分辨率的裁剪。在稀疏标注网格上、在小目标或严重遮挡的框上，这些条件往往同时不成立。因此本文把研究限定在只有检测的证据上，把外观视为自然的延伸而不是缺失的组件。我们实际做的对照，是把已发表跟踪器 BoostTrack \cite{stanjoevic2024boosttrack} 的关联评分移植进同一宿主，使只有评分组件发生变化。

\subsection{选择与风险控制}

给提案设门控是一个选择问题：阈值由数据选出，被选中阈值的实测风险天然偏乐观。Learn-then-Test \cite{angelopoulos2021learn} 在一个策略族上校准这一选择并给出有效的风险水平。相近的工作为跟踪中的候选转移集合给出覆盖保证 \cite{li2022pac}，或让检测框的不确定性在跟踪流程中传播 \cite{su2024motcup}。我们借用这条线的问题词汇来描述选点规则在做什么，并报告被选策略的实测失败比例，而不声称得到了校准的界。本文与这条线的区别不在统计，而在被测量的对象：闭环回放之后真正提交的重连动作，与它们产生的缺口恢复分开计数。
